\section{ESP32 Firmware Updates}

The existing STAR firmware requires additions for TCP communication as a fallback to SPI.

\subsection{TCP Server Module}

\begin{lstlisting}[language=C, caption=lib/star\_module\_tcp/include/star\_tcp\_server.h]
/**
 * @file star_tcp_server.h
 * @brief TCP server for command reception from RPi5
 *
 * Provides fallback communication when SPI is unavailable.
 */

#ifndef STAR_TCP_SERVER_H
#define STAR_TCP_SERVER_H

#include "esp_err.h"

#ifdef __cplusplus
extern "C" {
#endif

#define TCP_SERVER_PORT 5000
#define TCP_MAX_CLIENTS 2
#define TCP_RX_BUFFER_SIZE 256
#define TCP_TX_BUFFER_SIZE 256

/**
 * @brief Motor command callback type
 */
typedef struct {
    float left_velocity;
    float right_velocity;
    bool emergency_stop;
    uint64_t timestamp;
} tcp_motor_command_t;

typedef void (*tcp_command_callback_t)(const tcp_motor_command_t* command);

/**
 * @brief TCP server configuration
 */
typedef struct {
    uint16_t port;
    tcp_command_callback_t on_command;
    void* user_context;
} tcp_server_config_t;

/**
 * @brief TCP server handle
 */
typedef struct tcp_server_handle* tcp_server_handle_t;

/**
 * @brief Initialize and start TCP server
 *
 * @param config Server configuration
 * @param out_handle Output handle
 * @return esp_err_t ESP_OK on success
 */
esp_err_t star_tcp_server_init(
    const tcp_server_config_t* config,
    tcp_server_handle_t* out_handle
);

/**
 * @brief Stop and cleanup TCP server
 */
esp_err_t star_tcp_server_deinit(tcp_server_handle_t handle);

/**
 * @brief Get number of connected clients
 */
int star_tcp_server_get_client_count(tcp_server_handle_t handle);

#ifdef __cplusplus
}
#endif

#endif // STAR_TCP_SERVER_H
\end{lstlisting}

\begin{lstlisting}[language=C, caption=lib/star\_module\_tcp/src/star\_tcp\_server.c]
/**
 * @file star_tcp_server.c
 * @brief TCP server implementation
 */

#include "star_tcp_server.h"
#include "esp_log.h"
#include "esp_netif.h"
#include "lwip/sockets.h"
#include "cJSON.h"
#include "freertos/FreeRTOS.h"
#include "freertos/task.h"

static const char* TAG = "tcp_server";

struct tcp_server_handle {
    int listen_socket;
    tcp_command_callback_t callback;
    void* user_context;
    TaskHandle_t server_task;
    bool running;
    int client_count;
};

/**
 * @brief Parse JSON command from client
 */
static esp_err_t parse_json_command(
    const char* json_str,
    tcp_motor_command_t* out_cmd
) {
    cJSON* root = cJSON_Parse(json_str);
    if (root == NULL) {
        ESP_LOGE(TAG, "Failed to parse JSON");
        return ESP_ERR_INVALID_ARG;
    }

    memset(out_cmd, 0, sizeof(tcp_motor_command_t));

    cJSON* cmd = cJSON_GetObjectItem(root, "command");
    if (cmd && cJSON_IsString(cmd)) {
        if (strcmp(cmd->valuestring, "velocity") == 0) {
            cJSON* params = cJSON_GetObjectItem(root, "params");
            if (params) {
                cJSON* left = cJSON_GetObjectItem(params, "left");
                cJSON* right = cJSON_GetObjectItem(params, "right");

                if (left && cJSON_IsNumber(left)) {
                    out_cmd->left_velocity = (float)left->valuedouble;
                }
                if (right && cJSON_IsNumber(right)) {
                    out_cmd->right_velocity = (float)right->valuedouble;
                }
            }
        } else if (strcmp(cmd->valuestring, "emergency_stop") == 0) {
            out_cmd->emergency_stop = true;
        }
    }

    cJSON* ts = cJSON_GetObjectItem(root, "timestamp");
    if (ts && cJSON_IsNumber(ts)) {
        out_cmd->timestamp = (uint64_t)ts->valuedouble;
    }

    cJSON_Delete(root);
    return ESP_OK;
}

/**
 * @brief Build JSON response
 */
static void build_json_response(
    char* buffer,
    size_t size,
    bool success,
    const char* message
) {
    snprintf(buffer, size,
        "{\"success\":%s,\"message\":\"%s\"}",
        success ? "true" : "false",
        message
    );
}

/**
 * @brief Handle single client connection
 */
static void handle_client(
    tcp_server_handle_t handle,
    int client_socket
) {
    char rx_buffer[TCP_RX_BUFFER_SIZE];
    char tx_buffer[TCP_TX_BUFFER_SIZE];

    int len = recv(client_socket, rx_buffer, sizeof(rx_buffer) - 1, 0);

    if (len > 0) {
        rx_buffer[len] = '\0';
        ESP_LOGD(TAG, "Received: %s", rx_buffer);

        tcp_motor_command_t cmd;
        esp_err_t err = parse_json_command(rx_buffer, &cmd);

        if (err == ESP_OK && handle->callback) {
            handle->callback(&cmd);
            build_json_response(tx_buffer, sizeof(tx_buffer), true, "OK");
        } else {
            build_json_response(tx_buffer, sizeof(tx_buffer), false, "Parse error");
        }

        send(client_socket, tx_buffer, strlen(tx_buffer), 0);
    }
}

/**
 * @brief TCP server task
 */
static void tcp_server_task(void* pvParameters) {
    tcp_server_handle_t handle = (tcp_server_handle_t)pvParameters;

    struct sockaddr_in server_addr = {
        .sin_family = AF_INET,
        .sin_port = htons(TCP_SERVER_PORT),
        .sin_addr.s_addr = htonl(INADDR_ANY)
    };

    handle->listen_socket = socket(AF_INET, SOCK_STREAM, IPPROTO_TCP);
    if (handle->listen_socket < 0) {
        ESP_LOGE(TAG, "Failed to create socket");
        vTaskDelete(NULL);
        return;
    }

    int opt = 1;
    setsockopt(handle->listen_socket, SOL_SOCKET, SO_REUSEADDR, &opt, sizeof(opt));

    if (bind(handle->listen_socket, (struct sockaddr*)&server_addr,
             sizeof(server_addr)) < 0) {
        ESP_LOGE(TAG, "Failed to bind socket");
        close(handle->listen_socket);
        vTaskDelete(NULL);
        return;
    }

    if (listen(handle->listen_socket, TCP_MAX_CLIENTS) < 0) {
        ESP_LOGE(TAG, "Failed to listen");
        close(handle->listen_socket);
        vTaskDelete(NULL);
        return;
    }

    ESP_LOGI(TAG, "TCP server listening on port %d", TCP_SERVER_PORT);

    while (handle->running) {
        struct sockaddr_in client_addr;
        socklen_t addr_len = sizeof(client_addr);

        int client_socket = accept(
            handle->listen_socket,
            (struct sockaddr*)&client_addr,
            &addr_len
        );

        if (client_socket >= 0) {
            handle->client_count++;
            ESP_LOGI(TAG, "Client connected from %s",
                inet_ntoa(client_addr.sin_addr));

            handle_client(handle, client_socket);

            close(client_socket);
            handle->client_count--;
        }
    }

    close(handle->listen_socket);
    vTaskDelete(NULL);
}

esp_err_t star_tcp_server_init(
    const tcp_server_config_t* config,
    tcp_server_handle_t* out_handle
) {
    if (config == NULL || out_handle == NULL) {
        return ESP_ERR_INVALID_ARG;
    }

    tcp_server_handle_t handle = calloc(1, sizeof(struct tcp_server_handle));
    if (handle == NULL) {
        return ESP_ERR_NO_MEM;
    }

    handle->callback = config->on_command;
    handle->user_context = config->user_context;
    handle->running = true;
    handle->client_count = 0;

    BaseType_t ret = xTaskCreate(
        tcp_server_task,
        "tcp_server",
        4096,
        handle,
        5,
        &handle->server_task
    );

    if (ret != pdPASS) {
        free(handle);
        return ESP_FAIL;
    }

    *out_handle = handle;
    return ESP_OK;
}

esp_err_t star_tcp_server_deinit(tcp_server_handle_t handle) {
    if (handle == NULL) {
        return ESP_ERR_INVALID_ARG;
    }

    handle->running = false;

    // Wake up the accept() call
    if (handle->listen_socket >= 0) {
        shutdown(handle->listen_socket, SHUT_RDWR);
    }

    // Wait for task to finish
    vTaskDelay(pdMS_TO_TICKS(100));

    free(handle);
    return ESP_OK;
}

int star_tcp_server_get_client_count(tcp_server_handle_t handle) {
    return handle ? handle->client_count : 0;
}
\end{lstlisting}

\subsection{Main Application Integration}

Add TCP server to the main application:

\begin{lstlisting}[language=C, caption=Diff for ESP32-Firmware/src/main.c]
// Add to includes
#include "star_tcp_server.h"

// Add global handle
static tcp_server_handle_t g_tcp_server = NULL;

// Add command callback
static void on_tcp_command(const tcp_motor_command_t* cmd) {
    if (cmd->emergency_stop) {
        // Execute emergency stop
        motor_control_emergency_stop();
    } else {
        // Set motor velocities
        motor_control_set_velocity(cmd->left_velocity, cmd->right_velocity);
    }
}

// Add to app_main() after WiFi is connected
void app_main(void) {
    // ... existing initialization ...

    // Initialize TCP server for RPi5 communication
    tcp_server_config_t tcp_config = {
        .port = TCP_SERVER_PORT,
        .on_command = on_tcp_command,
        .user_context = NULL
    };

    esp_err_t err = star_tcp_server_init(&tcp_config, &g_tcp_server);
    if (err == ESP_OK) {
        ESP_LOGI(TAG, "TCP server started on port %d", TCP_SERVER_PORT);
    } else {
        ESP_LOGE(TAG, "Failed to start TCP server");
    }

    // ... rest of initialization ...
}
\end{lstlisting}

